% Generated draft: requires mathematical and source review.
\documentclass[11pt]{article}
\usepackage[margin=25mm]{geometry}
\usepackage{fontspec}
\setmainfont{Latin Modern Roman}
\usepackage{amsmath,amssymb,mathtools,unicode-math}
\setmathfont{Latin Modern Math}
\usepackage{enumitem,needspace,xurl,hyperref}
\setlength{\parindent}{0pt}
\setlength{\parskip}{5pt}
\setlength{\emergencystretch}{3em}
\newcommand{\N}{\mathbb{N}}
\newcommand{\Z}{\mathbb{Z}}
\newcommand{\Q}{\mathbb{Q}}
\newcommand{\R}{\mathbb{R}}
\newcommand{\C}{\mathbb{C}}
\newcommand{\F}{\mathbb{F}}
\newcommand{\A}{\mathbb{A}}
\newcommand{\PP}{\mathbb{P}}
\newcommand{\E}{\mathbb{E}}
\newcommand{\eps}{\varepsilon}
\newcommand{\dd}{\,\mathrm{d}}
\begin{document}
\section*{UnsolvedMath ID: 11000140}
\textbf{Source problem code:} catalogue entry
\subsection{Definitions and mathematical statement}
Let $\overline{\mathcal M}(F)$ be the compactified moduli space associated with the surface $F$ in the source. Problem 10 asks for a matrix model using the source's screens that calculates its (virtual) Euler characteristics. It explicitly says those invariants are already known, citing [28,56]; the requested novelty is the screen-based matrix-model calculation, not merely another numerical computation by unrelated means.

The supplied excerpt does not define screens or specify the compactification. It therefore does not justify silently identifying $\overline{\mathcal M}(F)$ with a chosen coarse or Deligne--Mumford space. Nor does it specify the matrix ensemble, potential, formal versus convergent integral, or the coefficients and weights by which the desired invariants should be recovered.

For context, in a suitable finite orbifold cell decomposition a weighted Euler sum has the form
\[
 \sum_{\text{cells }\sigma}\frac{(-1)^{\dim\sigma}}{|G_\sigma|},
\]
where $G_\sigma$ is the finite stabilizer. This illustrates the importance of automorphism weights, not a supplied screen formula or an identification of every possible meaning of virtual Euler characteristic. The requested model must specify its invariant convention and show that its prescribed expansion computes that invariant for the intended family of surfaces.
\subsection{Short English statement}
Build a matrix model based on screens that reproduces the already-known (virtual) Euler characteristics of the source's compactified moduli spaces. Screens, compactification and extraction weights are missing.
\subsection{Sources}
\begin{itemize}
\item Local registry: \url{data/unsolvedmath/problems.json}, permanent UnsolvedMath ID $11000140$, exact source code \texttt{catalogue entry}.
\item Farb, \emph{Problems on Mapping Class Groups and Related Topics}, Problem 10, PDF page 121: \url{https://www.math.uchicago.edu/~farb/papers/mcgbook.pdf}. Supplied access date: 2026-07-29; extraction: pdf-font-structure; repeated local numbering disambiguated by page among 329 headings.
\item Supplied references [28,56] for the known Euler characteristics lack bibliography entries here. The registry mentions Harer--Zagier and describes the screen-based approach as a Penner-style program, without exact references verifying the proposed model's current status.
\end{itemize}
\end{document}
